
          %%%%% ~~~~~~~~~~~~~~~~~~~~ %%%%%

\section{Estimation of
\;$\mu_{k} \;:=\, E_{Y}\!\left[\,Y_{k}\left\vert\,\Xobs\right.\right]$\;
and
\;$\sigma_{k}^{2} \;:=\, \Var_{Y}\!\left[\,Y_{k}\left\vert\,\Xobs\right.\right]$
}

We presented in Section \ref{sectionSEVANITheory} the underlying theory
of the SEVANI framework \cite{Beaumont2011} in a series of lemmas.
\vskip 0.2cm
\noindent
In particular, Lemma \ref{DecompositionMSETtilde} states that the the mean squared error
of a composite linear imputation estimator can be decomposed into the sum of three components:
the sampling variance component, the nonresponse variance component, and
the sampling-nonresponse covariance.
In turn, Lemma \ref{SamplingVariance}, Lemma \ref{NonresponseVariance}, and Lemma \ref{MixedVariance}
give respectively estimators for these three variance components.
\vskip 0.2cm
\noindent
The quantities
\;$\mu_{k} := E_{Y}\!\left[\,Y_{k}\left\vert\,\Xobs\right.\right]$\;
and
\;$\sigma_{k}^{2} := \Var_{Y}\!\left[\,Y_{k}\left\vert\,\Xobs\right.\right]$,\;
for \,$k \in s = s_{R} \sqcup s_{M}$,\,
appear in the formulas in
Lemma \ref{SamplingVariance}, Lemma \ref{NonresponseVariance}, and Lemma \ref{MixedVariance}.
The values of these quantities are unknown in practice, and must be estimated;
however, the SEVANI framework as described in \cite{Beaumont2011} does not specify how
\,$\mu_{k}$\, and \,$\sigma_{k}^{2}$\,
are to be estimated.
Instead the SEVANI framework \cite{Beaumont2011} treats the estimates of
\,$\mu_{k}$\, and \,$\sigma_{k}^{2}$\,
as input parameters of the framework.
\vskip 0.2cm
\noindent
On the other hand, the Methodology Guide of the SAS implementation of SEVANI does
contain discussions on the computations of
\,$\widehat{\mu}_{k}$\, and \,$\widehat{\sigma}_{k}^{2}$.\,
In this appendix, we give a summary of how these quantities are computed
by the SAS implementation of SEVANI, according to its Methodology Guide.

          %%%%% ~~~~~~~~~~~~~~~~~~~~ %%%%%

\vskip 0.5cm
\noindent
\textbf{Computations of
\,$\widehat{\mu}_{k}$\,
and
\,$\widehat{\sigma}_{k}^{2}$\,
in the SAS implementation of SEVANI}
\begin{itemize}
\item
    Prior to invoking SEVANI:
    \vskip 0.01cm
    \noindent
    The input sample data set should have undergone the imputation procedure by BANFF.
    The user specifies to BANFF the collection of imputation methods to apply and
    the order in which these methods are to be applied.
    During the BANFF imputation procedure, the selected sample
    \,$s = s_{R} \sqcup s_{M} \subset U$\,
    is divided into \,$J$\, ordered imputation classes.
\item
    Estimation of
    \;$\mu_{k} := E_{Y}\!\left[\,Y_{k}\left\vert\,\Xobs\right.\right]$\;
    and
    \;$\sigma_{k}^{2} := \Var_{Y}\!\left[\,Y_{k}\left\vert\,\Xobs\right.\right]$,\;
    for \,$k \in s = s_{R} \sqcup s_{M}$,\,
    is performed one imputation class at a time, and the order
    in which the estimation is performed follows that of the 
    imputation procedure upstream of SEVANI.
\item
    For each given imputation class (equivalently, at each modelling step),
    the collection of \textit{modelling contributors} is determined.
    The collection of modelling contributors will be used as training data
    for the given imputation class.
    \vskip 0.01cm
    \noindent
    The collection of modelling contributors is first initialized to the subset of respondents
    within the corresponding collection of \textit{imputation contributors}.
    (Recall that, for each imputation step, the collection of imputation contributors
    may contain nonrespondents for which the target variable has been imputed in preceding imputation steps.)
    If the collection of modelling contributors is not sufficiently large
    (to be used as training data), then it is enlarged by appropriately collapsing suitable imputation classes.
    \vskip 0.01cm
    \noindent
    Once the collection of modelling contributors has been determined,
    a predictive model
    -- based on the user-specified prediction technique for this imputation class/modelling step --
    is fitted to the \textit{adjusted data} (see next paragraph for more details)
    of the modelling contributors of this modelling step. 
    \vskip 0.01cm
    \noindent
    The SEVANI user may prescribe for each unit
    \,$k \in s = s_{R} \sqcup s_{M}$\,
    the following three modelling parameters:
    \begin{itemize}
    \item
        $\alpha_{k}$\, is the \textit{adjustment variable},
    \item
        $\nu_{k}$\, ``user-specified imputation model error variance'', and
    \item
        $\rho_{k}$\, is \textit{modelling weight}.
    \end{itemize}
    For each modelling contributor \,$k$,\, the \textit{adjusted target variable}
    \,$Y^{\prime}_{k}$\,
    is defined to be:
    \begin{equation*}
    Y^{\prime}_{k}
    \;\; := \;\;
        \dfrac{Y_{k}}{\alpha_{k}}
    \end{equation*}
    A predictive model is fitted to the following data:
    \begin{equation*}
    \left\{\;
        \left.
        (\,y^{\prime}_{k}\,,\,\mathbf{x}^{\textnormal{(obs)}}_{k}\,)
        \;\right\vert\;
        k \in \textnormal{Contrib}(C)
        \;\right\}
    \end{equation*}
\end{itemize}

          %%%%% ~~~~~~~~~~~~~~~~~~~~ %%%%%

\begin{equation*}
\left.\begin{array}{rcl}
Y_{k}
& = &
    \alpha_{k} \cdot Y^{\prime}_{k}
\\
\overset{{\color{white}1}}{E\!\left[\;
    Y^{\prime}_{k}
    \,\left\vert\;
    \mathbf{x}^{\textnormal{(obs)}}_{k}
    \right.
    \,\right]}
& = &
        \widehat{\mu}^{\prime}_{k}
\end{array}
\;\;\right\}
\textnormal{\large$\Longrightarrow$}\quad
E\!\left[\;
    \left.
    \overset{{\color{white}.}}{Y_{k}}
    \;\right\vert\;
    \mathbf{x}^{\textnormal{(obs)}}_{k}
    \,\right]
\;\; = \;\;
    \alpha_{k} \cdot \widehat{\mu}^{\prime}_{k}
\end{equation*}

% \begin{equation*}
% \Var\!\left[\;
%     Y^{\prime}_{k}
%     \;\left\vert\;
%     \mathbf{x}^{\textnormal{(obs)}}_{k}
%     \right.
%     \,\right]
% \;\; = \;\;
%     \Var\!\left[\;
%         \left.
%         \dfrac{Y_{k}}{\alpha_{k}}
%         \;\right\vert\;
%         \mathbf{x}^{\textnormal{(obs)}}_{k}
%         \,\right]
% \;\; = \;\;
%     \dfrac{\sigma_{k}^{2}}{\alpha_{k}^{2}}
% \;\; = \;\;
%     \left(\,\sigma^{\prime}_{C(k)}\,\right)^{2}
% \end{equation*}

\begin{equation*}
\left.\begin{array}{rcl}
Y_{k}
& = &
    \alpha_{k} \cdot Y^{\prime}_{k}
\\
\Var\!\left[\;
    Y^{\prime}_{k}
    \;\left\vert\;
    \mathbf{x}^{\textnormal{(obs)}}_{k}
    \right.
    \,\right]
& = &
    \overset{{\color{white}\textnormal{\large$1$}}}{
    \nu_{k}\cdot\left(\,\widehat{\sigma}^{\prime}_{C(k)}\,\right)^{2}
    }
\\
\left(\,\widehat{\sigma}^{\prime}_{C(k)}\,\right)^{2}
& = &
    \overset{{\color{white}\textnormal{\large$1$}}}{
    \dfrac{
        \underset{l\,\in\,C(k)}{\sum}\;
        \rho_{k}
        \left(\,y^{\prime}_{l} - \widehat{\mu}^{\prime}_{l}\,\right)^{2}
        }{
        \underset{l\,\in\,C(k)}{\sum}\;
        \rho_{l}\cdot\nu_{l}
        }
    }
\end{array}
\;\;\right\}
\Longrightarrow
\left\{\;\;
\begin{array}{ccl}
\Var\!\left[\;
    \left.
    \overset{{\color{white}.}}{Y_{k}}
    \;\right\vert\;
    \mathbf{x}^{\textnormal{(obs)}}_{k}
    \,\right]
& = &
    \alpha_{k}^{2}
    \cdot
    \nu_{k}
    \cdot 
    \left(\,\widehat{\sigma}^{\prime}_{C(k)}\,\right)^{2}
\\
& = &
    \overset{{\color{white}\textnormal{\small$1$}}}{
    \alpha_{k}^{2}
    \cdot
    \nu_{k}
    \cdot 
    \left(\,
    \dfrac{
        \underset{l\,\in\,C(k)}{\sum}\;
        \rho_{k}
        \left(\,y^{\prime}_{l} - \widehat{\mu}^{\prime}_{l}\,\right)^{2}
        }{
        \underset{l\,\in\,C(k)}{\sum}\;
        \rho_{l}\cdot\nu_{l}
        }
        \,\right)
    }
\end{array}
\right.
\end{equation*}
where
\begin{itemize}
\item
    $\alpha_{k}$\, is the \textit{adjustment variable},
\item
    $\nu_{k}$\, ``user-specified imputation model error variance'', and
\item
    $\rho_{k}$\, is \textit{modelling weight}.
\end{itemize}

          %%%%% ~~~~~~~~~~~~~~~~~~~~ %%%%%

\subsection{Parametric estimation: auxiliary value imputation}

\begin{equation*}
\begin{array}{ccccc}
\widehat{\mu}_{k}
& = &
    \widehat{E_{Y}}\!\left[\,Y_{k}\left\vert\,\Xobs\right.\right]
& = &
    \alpha_{k} \cdot x_{k}
\\ \\
\widehat{\sigma}_{k}^{2}
& = &
    \widehat{\Var_{Y}}\!\left[\,Y_{k}\left\vert\,\Xobs\right.\right]
& = &
    \alpha_{k}^{2} \cdot x_{k}
\end{array}
\end{equation*}
where
\,$\alpha_{k}$\, is the adjustment variable.

\begin{eqnarray*}
\widehat{\sigma}_{k}^{2}
& = &
    \widehat{\Var_{Y}}\!\left[\,Y_{k}\left\vert\,\Xobs\right.\right]
\\
& = &
    \alpha_{k}^{2}\cdot\widehat{\sigma}_{a,k}^{2}
\\
& = &
    \alpha_{k}^{2}\cdot\nu_{k}\cdot\widehat{\sigma}_{c}^{2}
\\
& = &
    \alpha_{k}^{2}\cdot\nu_{k}\cdot
    \dfrac{
        \underset{l\,\in\,C(k)}{\sum}\;
        \omega^{\textnormal{MOD}}_{l}
        \left(\,\dfrac{y_{l}}{\alpha_{l}} - \dfrac{\widehat{\mu}_{l}}{\alpha_{l}}\,\right)^{2}
        }{
        \underset{l\,\in\,C(k))}{\sum}\;
        \omega^{\textnormal{MOD}}_{l}\cdot\nu_{l}
        }
\end{eqnarray*}

          %%%%% ~~~~~~~~~~~~~~~~~~~~ %%%%%
